\documentclass{article}
\usepackage{graphicx} % Required for inserting images
\usepackage{algorithm} %
\usepackage{algpseudocode}
\usepackage{amsfonts}

\usepackage{tcolorbox}
\newtcolorbox{greenbox}[1]{colback=green!5!white,colframe=green!75!black,fonttitle=\bfseries,title=#1}


\usepackage{amsthm}
\usepackage{bm}

\newtheorem{definition}{(Definition)}
\newtheorem{theorem}{(Theorem)}
\newtheorem{lemma}{(Lemma)}
\newtheorem{corollary}{(Corollary)}

\newtheorem*{remark}{Remark}

\usepackage{amsmath}
\DeclareMathOperator*{\argmin}{argmin} % no space, limits underneath in displays



\setlength{\parindent}{0pt} % for getting rid of first line indentation


% ================================================
% =====================  BEGIN DOCUMENT 
% ================================================

\title{Stochastic Processes}
\author{Juan Pablo Bello Gonzalez}
\date{September 2024}

\begin{document}


\maketitle


\section{Lecture 1}
\textbf{Review of } $\mathbb{P}$, $\mathbb{E}$
, and $\int$

\textbf{A stochastic process} is a sequence of random variables indexed by set A, $(X_{\alpha})_{\alpha \in A}$. We can get creative with set A, in a time series it will be time but in poisson process it can take a '\textit{spatial}' meaning. The set A gives the sequence \textit{structure} through the mathematical structures in the set A such as ordering, closures, etc. 

Examples of stochastic processes: 
\begin{itemize}
\item Markov Chain (discrete and continuous)
\item Martingale 
\item Poisson process.
\item Time series. 
\item Continuous state Markov chains 
\item Empirical process
\end{itemize}


We want to formalize the assignment of probabilities to events in a sample space. 


We will now introduce $\sigma-$algebras as the  structure used to formalize events in a sample space; events to which we can assign a "probability" or "measure."

\subsection{Measure Theory}
Let $\Omega$ be our sample space. We are interested in the power set $\mathcal{P}(\Omega)$ which is the set of all subsets of $\Omega$. 

For instance, in an experiment in which we throw a three sided die, we have:
\begin{itemize}
\item $\Omega = \{1,2,3\}$
\item $P(\Omega) = \{ \emptyset, \{1\}, \{2\}, \{3\}, \{ 1,2 \}, \{ 1, 3\}, \{2,3\}, \{1,2,3\}  \}$
\end{itemize}

We are trying to formalize probability, so the idea of this \textit{first attempt} at it is to ask probability questions on $\Omega$ as questions regarding subsets of $\mathcal{P}(\Omega)$. For instance if we let X be the RV representing the side of the die, the event $(X=3)$ could be formalized as $X \in \{3\}$. This extends to more complex events such as $(X=1)\cup (X=2)$ translates to $X \in \{1,2\}$ .



 \textit{The power set is too general and too large}, so we restrict ourselves to the events given by a subset $\mathcal{A}$ of the powerset which satisfies some properties. This subset is called an \textbf{algebra}.
 
 \begin{definition}
 \textbf{Algebra}: A subset of the power set $\mathcal{A} \subset \mathcal{P}(\Omega)$ is an \textbf{algebra}, if: 
 \end{definition}
 \begin{enumerate}
 \item $\Omega \in \mathcal{A}$
 \item $A, B \in \mathcal{A} \implies A \cup B \in \mathbb{A}$
 \item $A \in \mathcal{A} \implies A^c = \Omega - A \in \mathcal{A}$
 \end{enumerate}
 
 \begin{corollary}
 If $\mathcal{A}$ is an algebra, then $A, B \in \mathcal{A} \implies A \cap B$.
 \end{corollary}
 \textit{Proof:} Using properties (2) and (3), $A,B \in \mathcal{A} \implies A^c, B^c \in \mathcal{A}$ so $A^c \cup B^c = (A \cap B)^c \in \mathcal{A}$ and since its compliment must also be in the algebra we have our result. 

  
  
  
\textbf{Questioning definitions:}
The power set is obviously an algebra. it contains the emptyset and the sample space set and all of its elements have their complements in it too. $\mathcal{P}$ is the largest algebra. So why do we bother with the definition of an algebra? because we are interested in the smaller algebras that live as subsets of $\mathcal{P}$ to have our theory work on the smaller conditional and other universes. \textbf{I THINK}. 







\end{document}
